\subsection{库默尔理论}

在本号中，我们用 n 表示一个整数 \textgreater0 ，并假设 \(\mu_{n}(\mathbf{K})\) 有 n 个元素；由 V，第78页，这也意味着 n 与 K 的特征指数互素，且 K 的代数闭包 \(\Omega\) 中的所有 n 次单位根都位于 K 中。

我们称 K 的扩张 L 是指数整除 n 的阿贝尔扩张，如果它是阿贝尔的（V，第77页，定义1）且其伽罗瓦群 \(\operatorname{Gal}(\mathrm{L}/\mathrm{K})\) 被 n 零化（V，第86页）。

设 A 是 \(K^{*}\) 的一个子集；我们用 \(\mathbf{K}(\mathbf{A}^{1/n})\) 表示 \(\Omega\) 的子扩张，它由所有 \(\theta \in \Omega\) （满足 \(\theta^{n} \in A\) ）所生成。

引理5. --- \(\mathbf{K}(\mathbf{A}^{1 / n})\) 是 \(\mathbf{K}\) 的一个指数整除 \(n\) 的阿贝尔扩张。

由于多项式 \(X^n - a\) （ \(a \in A\) ）在 \(K\) 上是可分的， \(L = K(A^{1/n})\) 是 \(K\) 的一个可分扩张，因而是伽罗瓦扩张。设 \(\sigma \in \operatorname{Gal}(L/K)\) ，并令 \(\theta \in \Omega\) 满足 \(\theta^n \in A\) 。我们有 \(\sigma(\theta)^n = \theta^n\) ；因此存在 \(\zeta \in \mu_n(\Omega) = \mu_n(K)\) 使得 \(\sigma(\theta) = \zeta\theta\) ；这蕴含 \(\sigma^n(\theta) = \zeta^n\theta = \theta\) ，从而 \(\sigma^n = 1\) 。若 \(\sigma'\) 是 \(\operatorname{Gal}(L/K)\) 的另一个元素，则存在 \(\zeta' \in \mu_n(K)\) 使得 \(\sigma'(\theta) = \zeta'\theta\) ，从而 \(\sigma' \sigma(\theta) = \zeta\zeta'\theta = \sigma\sigma'(\theta)\) ，因此 \(\sigma' \sigma = \sigma\sigma'\) 。

引理6. --- 设 L 是 K 的伽罗瓦扩张。存在唯一的映射 \((\sigma, a) \mapsto \langle \sigma, a \rangle\) （从 \(\operatorname{Gal}(\mathrm{L} / \mathrm{K}) \times ((\mathrm{L}^n \cap \mathrm{K}^*) / \mathrm{K}^{*n})\) 到 \(\mu_n(\mathbf{K})\) 中），使得对每个 \(\sigma \in \operatorname{Gal}(\mathrm{L} / \mathrm{K})\) 以及每个元素 \(\theta \in \mathrm{L}^*\) （满足 \(\theta^n \in \mathbf{K}\) ），用 \(\bar{\theta}^n\) 表示 \(\theta^n \bmod \mathrm{K}^{*n}\) 的剩余类时，我们有：

\[
\left\langle \sigma , \bar {\theta} ^ {n} \right\rangle = \sigma (\theta) / \theta .\tag{13}
\]

该映射是双乘性的。

实际上，(13)式右端是一个 n 次单位根，它只依赖于模 \(K^{*n}\) 的 \(\theta^{n}\) 剩余类；这证明了第一个断言。第二个断言不难验证。

对于 K 的每一个伽罗瓦扩张 L，我们记

\[
\begin{array}{l} k _ {\mathrm{L}} \colon (\mathrm{L} ^ {n} \cap \mathrm{K} ^ {*}) / \mathrm{K} ^ {* n} \to \operatorname{Hom} \left(\operatorname{Gal} (\mathrm{L} / \mathrm{K}),   \mu_ {n} (\mathrm{K})\right), \\ k _ {\mathrm{L}} ^ {\prime} \colon \operatorname{Gal} (\mathrm{L} / \mathrm{K}) \to \operatorname{Hom} \left((\mathrm{L} ^ {n} \cap \mathrm{K} ^ {*}) / \mathrm{K} ^ {* n},   \mu_ {n} (\mathrm{K})\right), \end{array}
\]

为由上述双乘性映射导出的同态（V，第87页）。

命题9. --- 对于 K 的每一个有限次伽罗瓦扩张 L，同态 \(k_{L}\) 是双射的。

设 \(\theta \in L^{*}\) 满足 \(\theta^n \in K\) ，且 \(\theta^n \bmod K^{*n}\) 的剩余类属于 \(k_{L}\) 的核。对每个 \(\sigma \in \operatorname{Gal}(L/K)\) ，根据定义我们有 \(\sigma(\theta) = \theta\) ；因此 \(\theta \in K^{*}\) 且 \(\theta^n \in K^{*n}\) 。这证明了 \(k_{L}\) 的单射性。现在设 \(f: \operatorname{Gal}(L/K) \to \mu_n(K)\) 为一个同态；对所有 \(\sigma, \tau \in \operatorname{Gal}(L/K)\) 我们有

\[
f (\sigma \tau) = f (\sigma) f (\tau) = f (\sigma). \sigma f (\tau), \quad f (\sigma) ^ {n} = 1.
\]

由 V，第65页，推论1，存在 \(\theta \in L^{*}\) 使得 \(f(\sigma) = \sigma(\theta)/\theta\) （对所有 \(\sigma \in \operatorname{Gal}(L/K)\) ）；由于 \(f(\sigma)^{n} = 1\) ，我们有 \(\sigma(\theta^{n}) = \theta^{n}\) （对所有 \(\sigma \in \operatorname{Gal}(L/K)\) ），因此 \(\theta^{n} \in K^{*}\) ；若 \(a\) 是 \(\theta^{n} \bmod K^{*n}\) 的剩余类，则根据定义我们有 \(f(\sigma) = \langle \sigma, a \rangle\) （对 \(\sigma \in \operatorname{Gal}(L/K)\) ），因此 \(f = k_{L}(a)\) 。

推论. --- 若 L 是 K 的伽罗瓦扩张，则同态 \(k_{L}\) 是单射的，且其像是群 \(\operatorname{Hom}_{c}(\operatorname{Gal}(\operatorname{L}/\operatorname{K}), \mu_{n}(\operatorname{K}))\) （由拓扑群 \(\operatorname{Gal}(\operatorname{L}/\operatorname{K})\) 到离散群 \(\mu_{n}(\operatorname{K})\) 中的连续同态组成的群）。

这由前述内容立即得出，利用下述事实：L 是有限次伽罗瓦扩张 \(L_{i}\) 的递增有向并，且 \(\operatorname{Gal}(L/K)\) 到 \(\mu_{n}(K)\) 中的同态是连续的，当且仅当它可以通过 \(\operatorname{Gal}(L_{i}/K)\) （ \(\operatorname{Gal}(L/K)\) 的某个商群）分解。

定理4. --- a) 映射 \(\mathrm{H} \mapsto \mathrm{K}\left(\mathrm{H}^{1/n}\right)\) 是 \(\mathbf{K}^*\) 的包含 \(\mathbf{K}^{*n}\) 的子群组成的集合到指数整除 \(n\) 的、 \(\Omega\) 的阿贝尔子扩张组成的集合上的一个保持包含关系的双射。逆映射为 \(\mathrm{L} \mapsto \mathrm{L}^n \cap \mathbf{K}^*\) 。

\begin{enumerate}
\def\labelenumi{\alph{enumi})}
\setcounter{enumi}{1}
\tightlist
\item
  对 \(\mathbf{K}^*\) 的每个包含 \(\mathbf{K}^{*n}\) 的子群 H，同态
\end{enumerate}

\[
k ^ {\prime}: \operatorname{Gal} \left(\mathrm{K} \left(\mathrm{H} ^ {1 / n}\right) / \mathrm{K}\right)\rightarrow \operatorname{Hom} \left(\mathrm{H} / \mathrm{K} ^ {* n}, \mu_ {n} (\mathrm{K})\right)
\]

是双射的，且当群 \(\operatorname{Hom}(\mathrm{H} / \mathrm{K}^{*n},\mu_n(\mathrm{K}))\) 赋以简单收敛拓扑时，它是一个同胚。

\begin{enumerate}
\def\labelenumi{\alph{enumi})}
\setcounter{enumi}{2}
\tightlist
\item
  设 H 是 \(K^{*}\) 的一个包含 \(K^{*n}\) 的子群。对每个 \(a \in H/K^{*n}\) ，令 \(\theta_{a}\) 为 \(\Omega\) 的一个元素，使得 \(\theta_{a}^{n}\) 是 a 在 H 中的一个代表元。那么诸 \(\theta_{a}\) （ \(a \in H/K^{*n}\) ）构成 K-向量空间 \(\mathrm{K}(H^{1/n})\) 的一组基。特别地，
\end{enumerate}

\[
[ \mathrm{K} (\mathrm{H} ^ {1 / n}): \mathrm{K} ] = (\mathrm{H}: \mathrm{K} ^ {* n}).
\]

\begin{enumerate}
\def\labelenumi{\Alph{enumi})}
\tightlist
\item
  对 K 的每个指数整除 n 的阿贝尔扩张 L，令 \(H_{L} = L^{n} \cap K^{*}\) 。若 \([L : K]\) 是有限的，则同态 \(k_{L}'\) （从 \(\operatorname{Gal}(L/K)\) 到 \(\operatorname{Hom}(H_{L}, \mu_{n}(K))\) 中）由命题9以及 V，第88页，命题8可知是双射的。
\end{enumerate}

由于 K 的每个指数整除 n 的阿贝尔扩张都是指数整除 n 的有限次阿贝尔子扩张的递增有向并，我们通过取逆极限推出：对 K 的每个指数整除 n 的阿贝尔扩张 L， \(k_{L}^{\prime}\) 是拓扑群的一个同胚。

\begin{enumerate}
\def\labelenumi{\Alph{enumi})}
\setcounter{enumi}{1}
\item
  设 L 是 K 的有限次阿贝尔扩张，其指数整除 n，并设 \(L' = \mathbf{K}(\mathbf{H}_{L}'^{1/n})\) ；这是 L 的一个子扩张；此外 \(H_{L'}\) 包含 \(H_{L}\) ，因此等于它。由于同态 \(k_{L}'\) 和 \(k_{L'}\) 由 A) 以及 \(H_{L} = H_{L'}\) 可知是双射的，群 \(\operatorname{Gal}(L/K)\) 和 \(\operatorname{Gal}(L'/K)\) 具有相同的阶，因此相等。这表明 \(L' = L\) ，从而 \(L = \mathbf{K}(\mathbf{H}_{L}'^{1/n})\) 。若 L 是 K 的指数整除 n 的阿贝尔扩张，则我们有 \(\mathbf{K}(\mathbf{H}_{L}'^{1/n}) = \mathbf{L}\) ，因为 \(\mathbf{K}(\mathbf{H}_{L}'^{1/n})\) 是 L 的一个子扩张，且包含 L 的每一个有限次子扩张。
\item
  设 H 是 \(K^{*}\) 的一个包含 \(K^{*n}\) 的子群；令 \(L = K(H^{1/n})\) ，则这是 K 的一个指数整除 n 的阿贝尔扩张（V，第88页，引理5）。我们有 \(H \subset H_{L}\) ，由此我们得到一个被 n 零化的交换群的正合序列
\end{enumerate}

\[
\{1 \} \rightarrow \mathrm{H} / \mathrm{K} ^ {* n} \rightarrow \mathrm{H} _ {\mathrm{L}} / \mathrm{K} ^ {* n} \rightarrow \mathrm{H} _ {\mathrm{L}} / \mathrm{H} \rightarrow \{1 \}.
\]

由此我们得到一个精确序列

\[
\left\{1 \right\}\rightarrow \mathrm{Hom} \left(\mathrm{H} _ {\mathrm{L}} / \mathrm{H}, \mu_ {n} (\mathrm{K})\right)\rightarrow \mathrm{Hom} \left(\mathrm{H} _ {\mathrm{L}} / \mathrm{K} ^ {* n}, \mu_ {n} (\mathrm{K})\right) \stackrel {u} {\rightarrow} \mathrm{Hom} \left(\mathrm{H} / \mathrm{K} ^ {* n}, \mu_ {n} (\mathrm{K})\right)
\]

其中 \(u\) 是限制同态。

如果我们把 \(\operatorname{Hom}(\mathrm{H}_{\mathrm{L}}/\mathrm{K}^{*n},\mu_{n}(\mathrm{K}))\) 与 \(\operatorname{Gal}(\mathrm{L}/\mathrm{K})\) （借助同构 \(k_{L}^{\prime}\) ）等同起来，则 u 的核等同于由所有 \(\sigma\in\operatorname{Gal}(\mathrm{L}/\mathrm{K})\) （满足 \(\sigma(\theta)=\theta\) ，对所有 \(\theta\in H^{1/n}\) ）组成的集合。由此可知 u 是单射，因此 \(\operatorname{Hom}(\mathrm{H}_{\mathrm{L}}/\mathrm{H},\mu_{n}(\mathrm{K}))\) 约化为 \(\{1\}\) ；由 V，第87页，推论1，我们从而有 \(H=H_{L}\) 。这完成了 a) 和 b) 的证明。

\begin{enumerate}
\def\labelenumi{\Alph{enumi})}
\setcounter{enumi}{3}
\tightlist
\item
  现在证明 c)。若 \(a, b \in H\) ，我们有 \(\theta_{a}\theta_{b}/\theta_{ab} \in K\) 。由此可知， \(\mathrm{K}(\mathrm{H}^{1/n})\) 的由诸 \(\theta_{a}\) 生成的向量子空间在乘法下稳定，因此与 \(\mathrm{K}(\mathrm{H}^{1/n})\) 重合。剩下的只需证明诸 \(\theta_{a}\) 线性无关；为此我们显然可以假设 \(H/K^{*n}\) 是有限的；于是由 b) 及 V，第87页，推论2，得 \([\mathrm{K}(\mathrm{H}^{1/n}): \mathrm{K}] = (\mathrm{Gal}(\mathrm{K}(\mathrm{H}^{1/n})/\mathrm{K}) : \{1\}) = (\mathrm{H}: \mathrm{K}^{*n})\) ；由于 \(\theta_{a}\) 的个数等于 \((\mathrm{H}: \mathrm{K}^{*n})\) 且它们生成 K-向量空间 \(\mathrm{K}(\mathrm{H}^{1/n})\) ，所以它们线性无关。
\end{enumerate}

例子. --- 1) 存在 K 的指数整除 n 的最大阿贝尔扩张，包含于 \(\Omega\) ；它通过将 K 的所有元素的 n 次根添加到 K 而得到；其伽罗瓦群可与 \(\operatorname{Hom}(\mathbf{K}^{*}/\mathbf{K}^{*n}, \mu_{n}(\mathbf{K}))\) 等同，因此也与 \(\operatorname{Hom}(\mathbf{K}^{*}, \mu_{n}(\mathbf{K}))\) 等同。

* 2) 我们取 \(\mathbf{K} = \mathbf{Q}\) 和 \(n = 2\) 。则 \(\mathbf{Q}^{*}/\mathbf{Q}^{*2}\) 是一个 \(\mathbf{F}_{2}\) -向量空间，其基是 \(\{-1\}\) 与所有素数组成的集合的并。 \(\mathbf{Q}\) 的包含在 \(\mathbf{C}\) 中的指数为2的最大阿贝尔扩张因此是子域 \(\mathbf{Q}(i, \sqrt{2}, \sqrt{3}, \sqrt{5}, \ldots)\) （ \(\mathbf{C}\) 的一个子域）。其伽罗瓦群由将元素 \(i\) , \(\sqrt{2}\) , \(\sqrt{3}\) , \(\sqrt{5}\) 等分别乘以 \(\pm 1\) 所得到的所有自同构组成。*

\begin{enumerate}
\def\labelenumi{\arabic{enumi})}
\setcounter{enumi}{2}
\item
  设 \(L\) 是 \(K\) 的一个次数为 \(n\) 的循环扩张；则群 \((L^n \cap K^*) / K^{*n}\) 是阶为 \(n\) 的循环群。若 \(a \in K^*\) 使得 \(a \mod K^{*n}\) 的剩余类是该群的生成元，则 \(L\) 是 \(K\) -同构于 \(K[X] / (X^n - a)\) 的，且群 \(\operatorname{Gal}(L / K)\) 由 \(n\) 个自同构组成，它们将 \(X\) 映到 \(\zeta X\) （ \(\zeta \in \mu_n(K)\) ）。
\item
  反之，设 \(a \in \mathbf{K}^*\) ，并令 \(r\) 为最小整数 \(> 0\) （使得 \(a^r \in \mathbf{K}^{*n}\) ）；则子域 \(L\) （ \(\Omega\) 的、由多项式 \(X^n - a\) 的根生成的）是 \(K\) 的一个次数为 \(r\) 的循环扩张。特别地， \(X^n - a\) 不可约，当且仅当 \(r = n\) 。
\end{enumerate}

注记. --- 设 \(a \in K^{*}\) ，并设 r 为最小整数 \textgreater0 （使得 \(a^{r} \in K^{n}\) ）。设 B 为多项式 \(X^{n/r} - a\) 在 K 中的根组成的集合；于是我们有

\[
\mathbf {X} ^ {n} - a = \prod_ {b \in B} \left(\mathbf {X} ^ {r} - b\right),\tag{14}
\]

通过以 \(X^{r}\) 代替关系式 \(T^{n/r}-a=\Pi(T-b)\) 中的 T 而得。根据例4，每个多项式 \(X^{r}-b\) 都是不可约的，因此（14）是 \(X^{n}-a\) 在 K{[}X{]} 中的不可约多项式分解。

